
\documentclass{article}
\usepackage{colortbl}
\usepackage{makecell}
\usepackage{multirow}
\usepackage{supertabular}

\begin{document}

\newcounter{utterance}

\centering \large Interaction Transcript for game `codenames', experiment `association\_easy', episode 3 with deepseek{-}v4{-}pro.
\vspace{24pt}

{ \footnotesize  \setcounter{utterance}{1}
\setlength{\tabcolsep}{0pt}
\begin{supertabular}{c@{$\;$}|p{.15\linewidth}@{}p{.15\linewidth}p{.15\linewidth}p{.15\linewidth}p{.15\linewidth}p{.15\linewidth}}
    \# & $\;$A & \multicolumn{4}{c}{Game Master} & $\;\:$B\\
    \hline

    \theutterance \stepcounter{utterance}  
    & & & \multicolumn{2}{p{0.3\linewidth}}{
        \cellcolor[rgb]{0.9,0.9,0.9}{
            \makecell[{{p{\linewidth}}}]{
                \texttt{\tiny{[GM$|$GM]}}
                {'hidden': {'team': ['strzyżyk', 'huragan', 'oszustwo', 'rozbój', 'przekroczenie prędkości', 'deszcz', 'kardynał', 'tornado', 'tajfun'], 'innocent': ['stopa', 'koszula', 'jard', 'krawat', 'pasek', 'kilometr', 'bluza'], 'opponent': ['nagietek', 'róża', 'księgowy', 'pistolet', 'bąk', 'pelargonia', 'prawnik', 'mechanik'], 'assassin': ['cymbał']}, 'revealed': {'team': {'team': [], 'innocent': [], 'opponent': [], 'assassin': []}, 'opponent': {'team': [], 'innocent': [], 'opponent': [], 'assassin': []}}}
            }
        }
    }
    & & \\ \\

    \theutterance \stepcounter{utterance}  
    & & \multicolumn{4}{p{0.6\linewidth}}{
        \cellcolor[rgb]{0.9,0.9,0.9}{
            \makecell[{{p{\linewidth}}}]{
                \texttt{\tiny{[P1$\langle$GM]}}
                \texttt{Podaj jedną wskazówkę składającą się z pojedynczego słowa, która odnosi się do niektórych słów twojej drużyny z poniższej listy. Możesz wybrać, do ilu i do jakich słów odnosi się twoja wskazówka.} \\
\texttt{Celem jest, aby twój partner odgadł wszystkie słowa twojej drużyny na podstawie podanych wskazówek, zanim drużyna przeciwna odkryje wszystkie swoje. Jeśli twój partner odgadnie słowo 'przeciwnika' lub 'rozpraszające', jego tura zgadywania kończy się i drużyna przeciwna wykonuje turę. Jeśli twój partner odgadnie słowo 'zabójcze', gra kończy się i twoja drużyna natychmiast przegrywa.} \\
\texttt{Im więcej słów odnosi się do twojej wskazówki, tym szybciej możesz wygrać grę, ale tym trudniejsze mogą być skojarzenia do odgadnięcia dla twojego partnera, więc wybieraj słowo wskazówki mądrze. Słowo wskazówki musi być semantycznie powiązane ze słowami docelowymi, nie może być jednym ze słów na listach ani zawierać ich części.} \\
\\ 
\texttt{Zawsze podawaj swoją wskazówkę składającą się z pojedynczego słowa i oddzieloną przecinkami listę powiązanych słów docelowych w następującym formacie i staraj się, aby twoje odpowiedzi były jak najkrótsze, nigdy nie dołączaj żadnego innego tekstu niż jest to wymagane w tym formacie:} \\
\\ 
\texttt{WSKAZÓWKA: <SŁOWO>} \\
\texttt{CELE: <CELE>} \\
\\ 
\texttt{Słowa twojej drużyny to: strzyżyk, huragan, oszustwo, rozbój, przekroczenie prędkości, deszcz, kardynał, tornado, tajfun.} \\
\texttt{Słowa drużyny przeciwnika to: nagietek, róża, księgowy, pistolet, bąk, pelargonia, prawnik, mechanik.} \\
\texttt{Słowa rozpraszające to: stopa, koszula, jard, krawat, pasek, kilometr, bluza.} \\
\texttt{Słowa zabójcze to: cymbał.} \\
            }
        }
    }
    & \\ \\

    \theutterance \stepcounter{utterance}  
    & \multicolumn{4}{p{0.6\linewidth}}{
        \cellcolor[rgb]{0.9,0.9,0.9}{
            \makecell[{{p{\linewidth}}}]{
                \texttt{\tiny{[P1$\rangle$GM]}}
                \texttt{WSKAZÓWKA: wiatr} \\
\texttt{CELE: huragan, tornado, tajfun} \\
            }
        }
    }
    & & \\ \\

    \theutterance \stepcounter{utterance}  
    & & & \multicolumn{2}{p{0.3\linewidth}}{
        \cellcolor[rgb]{0.9,0.9,0.9}{
            \makecell[{{p{\linewidth}}}]{
                \texttt{\tiny{[GM$|$GM]}}
                \texttt{wiatr} \\
            }
        }
    }
    & & \\ \\

    \theutterance \stepcounter{utterance}  
    & & & \multicolumn{2}{p{0.3\linewidth}}{
        \cellcolor[rgb]{0.9,0.9,0.9}{
            \makecell[{{p{\linewidth}}}]{
                \texttt{\tiny{[GM$|$GM]}}
                ['huragan', 'tornado', 'tajfun']
            }
        }
    }
    & & \\ \\

    \theutterance \stepcounter{utterance}  
    & & & \multicolumn{2}{p{0.3\linewidth}}{
        \cellcolor[rgb]{0.9,0.9,0.9}{
            \makecell[{{p{\linewidth}}}]{
                \texttt{\tiny{[GM$|$GM]}}
                {'word': 'huragan', 'assignment': 'team'}
            }
        }
    }
    & & \\ \\

    \theutterance \stepcounter{utterance}  
    & & & \multicolumn{2}{p{0.3\linewidth}}{
        \cellcolor[rgb]{0.9,0.9,0.9}{
            \makecell[{{p{\linewidth}}}]{
                \texttt{\tiny{[GM$|$GM]}}
                {'word': 'tornado', 'assignment': 'team'}
            }
        }
    }
    & & \\ \\

    \theutterance \stepcounter{utterance}  
    & & & \multicolumn{2}{p{0.3\linewidth}}{
        \cellcolor[rgb]{0.9,0.9,0.9}{
            \makecell[{{p{\linewidth}}}]{
                \texttt{\tiny{[GM$|$GM]}}
                {'word': 'tajfun', 'assignment': 'team'}
            }
        }
    }
    & & \\ \\

    \theutterance \stepcounter{utterance}  
    & & \multicolumn{4}{p{0.6\linewidth}}{
        \cellcolor[rgb]{0.9,0.9,0.9}{
            \makecell[{{p{\linewidth}}}]{
                \texttt{\tiny{[GM$\rangle$P2]}}
                \texttt{Podaj oddzieloną przecinkami listę maksymalnie 3 słów z poniższej listy, które najlepiej odnoszą się lub są najbardziej związane ze słowem 'wiatr'. Zawsze rozpoczynaj swoją listę odgadnięć od ZGADNIĘCIE:  i nie dołączaj żadnego innego tekstu w swojej odpowiedzi.} \\
\\ 
\texttt{mechanik, rozbój, deszcz, krawat, kardynał, strzyżyk, róża, stopa, nagietek, pistolet, bąk, prawnik, pelargonia, cymbał, kilometr, tajfun, tornado, pasek, bluza, księgowy, oszustwo, jard, przekroczenie prędkości, koszula, huragan} \\
            }
        }
    }
    & \\ \\

    \theutterance \stepcounter{utterance}  
    & & & \multicolumn{4}{p{0.6\linewidth}}{
        \cellcolor[rgb]{0.9,0.9,0.9}{
            \makecell[{{p{\linewidth}}}]{
                \texttt{\tiny{[GM$\langle$P2]}}
                \texttt{ZGADNIĘCIE: tajfun, tornado, huragan} \\
            }
        }
    }
     \\ \\

    \theutterance \stepcounter{utterance}  
    & & & \multicolumn{2}{p{0.3\linewidth}}{
        \cellcolor[rgb]{0.9,0.9,0.9}{
            \makecell[{{p{\linewidth}}}]{
                \texttt{\tiny{[GM$|$GM]}}
                ['tajfun', 'tornado', 'huragan']
            }
        }
    }
    & & \\ \\

    \theutterance \stepcounter{utterance}  
    & & & \multicolumn{2}{p{0.3\linewidth}}{
        \cellcolor[rgb]{0.9,0.9,0.9}{
            \makecell[{{p{\linewidth}}}]{
                \texttt{\tiny{[GM$|$GM]}}
                {'word': 'tajfun', 'assignment': 'team'}
            }
        }
    }
    & & \\ \\

    \theutterance \stepcounter{utterance}  
    & & & \multicolumn{2}{p{0.3\linewidth}}{
        \cellcolor[rgb]{0.9,0.9,0.9}{
            \makecell[{{p{\linewidth}}}]{
                \texttt{\tiny{[GM$|$GM]}}
                {'word': 'tajfun', 'assignment': 'team'}
            }
        }
    }
    & & \\ \\

    \theutterance \stepcounter{utterance}  
    & & & \multicolumn{2}{p{0.3\linewidth}}{
        \cellcolor[rgb]{0.9,0.9,0.9}{
            \makecell[{{p{\linewidth}}}]{
                \texttt{\tiny{[GM$|$GM]}}
                {'word': 'tornado', 'assignment': 'team'}
            }
        }
    }
    & & \\ \\

    \theutterance \stepcounter{utterance}  
    & & & \multicolumn{2}{p{0.3\linewidth}}{
        \cellcolor[rgb]{0.9,0.9,0.9}{
            \makecell[{{p{\linewidth}}}]{
                \texttt{\tiny{[GM$|$GM]}}
                {'word': 'tornado', 'assignment': 'team'}
            }
        }
    }
    & & \\ \\

    \theutterance \stepcounter{utterance}  
    & & & \multicolumn{2}{p{0.3\linewidth}}{
        \cellcolor[rgb]{0.9,0.9,0.9}{
            \makecell[{{p{\linewidth}}}]{
                \texttt{\tiny{[GM$|$GM]}}
                {'word': 'huragan', 'assignment': 'team'}
            }
        }
    }
    & & \\ \\

    \theutterance \stepcounter{utterance}  
    & & & \multicolumn{2}{p{0.3\linewidth}}{
        \cellcolor[rgb]{0.9,0.9,0.9}{
            \makecell[{{p{\linewidth}}}]{
                \texttt{\tiny{[GM$|$GM]}}
                {'word': 'huragan', 'assignment': 'team'}
            }
        }
    }
    & & \\ \\

    \theutterance \stepcounter{utterance}  
    & & & \multicolumn{2}{p{0.3\linewidth}}{
        \cellcolor[rgb]{0.9,0.9,0.9}{
            \makecell[{{p{\linewidth}}}]{
                \texttt{\tiny{[GM$|$GM]}}
                {'word': 'nagietek', 'assignment': 'opponent'}
            }
        }
    }
    & & \\ \\

    \theutterance \stepcounter{utterance}  
    & & & \multicolumn{2}{p{0.3\linewidth}}{
        \cellcolor[rgb]{0.9,0.9,0.9}{
            \makecell[{{p{\linewidth}}}]{
                \texttt{\tiny{[GM$|$GM]}}
                {'hidden': {'team': ['strzyżyk', 'oszustwo', 'rozbój', 'przekroczenie prędkości', 'deszcz', 'kardynał'], 'innocent': ['stopa', 'koszula', 'jard', 'krawat', 'pasek', 'kilometr', 'bluza'], 'opponent': ['róża', 'księgowy', 'pistolet', 'bąk', 'pelargonia', 'prawnik', 'mechanik'], 'assassin': ['cymbał']}, 'revealed': {'team': {'team': ['tajfun', 'tornado', 'huragan'], 'innocent': [], 'opponent': [], 'assassin': []}, 'opponent': {'team': [], 'innocent': [], 'opponent': ['nagietek'], 'assassin': []}}}
            }
        }
    }
    & & \\ \\

    \theutterance \stepcounter{utterance}  
    & & \multicolumn{4}{p{0.6\linewidth}}{
        \cellcolor[rgb]{0.9,0.9,0.9}{
            \makecell[{{p{\linewidth}}}]{
                \texttt{\tiny{[P1$\langle$GM]}}
                \texttt{Słowa tajfun, tornado, huragan zostały odgadnięte prawidłowo. Tura twojego kolegi z drużyny się tam skończyła.} \\
\texttt{Teraz podaj kolejną wskazówkę odnoszącą się do niektórych z pozostałych słów twojego zespołu oraz listę powiązanych słów docelowych. Pamiętaj, aby rozpocząć wskazówkę od 'WSKAZÓWKA: ', wstawić nową linię i rozpocząć listę słów docelowych oddzielonych przecinkami od 'CELE: '. Uwaga: niektóre słowa zostały usunięte z list w porównaniu do poprzednich żądań.} \\
\\ 
\texttt{Twoje pozostałe słowa zespołu to: strzyżyk, oszustwo, rozbój, przekroczenie prędkości, deszcz, kardynał.} \\
\texttt{Pozostałe słowa twojego przeciwnika to: nagietek, róża, księgowy, pistolet, bąk, pelargonia, prawnik, mechanik.} \\
\texttt{Pozostałe słowa rozpraszające to: stopa, koszula, jard, krawat, pasek, kilometr, bluza.} \\
\texttt{Pozostałe słowa zabójcy to: cymbał.} \\
            }
        }
    }
    & \\ \\

    \theutterance \stepcounter{utterance}  
    & \multicolumn{4}{p{0.6\linewidth}}{
        \cellcolor[rgb]{0.9,0.9,0.9}{
            \makecell[{{p{\linewidth}}}]{
                \texttt{\tiny{[P1$\rangle$GM]}}
                \texttt{WSKAZÓWKA: przestępstwo} \\
\texttt{CELE: oszustwo, rozbój, przekroczenie prędkości} \\
            }
        }
    }
    & & \\ \\

    \theutterance \stepcounter{utterance}  
    & & & \multicolumn{2}{p{0.3\linewidth}}{
        \cellcolor[rgb]{0.9,0.9,0.9}{
            \makecell[{{p{\linewidth}}}]{
                \texttt{\tiny{[GM$|$GM]}}
                \texttt{przestępstwo} \\
            }
        }
    }
    & & \\ \\

    \theutterance \stepcounter{utterance}  
    & & & \multicolumn{2}{p{0.3\linewidth}}{
        \cellcolor[rgb]{0.9,0.9,0.9}{
            \makecell[{{p{\linewidth}}}]{
                \texttt{\tiny{[GM$|$GM]}}
                ['oszustwo', 'rozbój', 'przekroczenie prędkości']
            }
        }
    }
    & & \\ \\

    \theutterance \stepcounter{utterance}  
    & & & \multicolumn{2}{p{0.3\linewidth}}{
        \cellcolor[rgb]{0.9,0.9,0.9}{
            \makecell[{{p{\linewidth}}}]{
                \texttt{\tiny{[GM$|$GM]}}
                {'word': 'oszustwo', 'assignment': 'team'}
            }
        }
    }
    & & \\ \\

    \theutterance \stepcounter{utterance}  
    & & & \multicolumn{2}{p{0.3\linewidth}}{
        \cellcolor[rgb]{0.9,0.9,0.9}{
            \makecell[{{p{\linewidth}}}]{
                \texttt{\tiny{[GM$|$GM]}}
                {'word': 'rozbój', 'assignment': 'team'}
            }
        }
    }
    & & \\ \\

    \theutterance \stepcounter{utterance}  
    & & & \multicolumn{2}{p{0.3\linewidth}}{
        \cellcolor[rgb]{0.9,0.9,0.9}{
            \makecell[{{p{\linewidth}}}]{
                \texttt{\tiny{[GM$|$GM]}}
                {'word': 'przekroczenie prędkości', 'assignment': 'team'}
            }
        }
    }
    & & \\ \\

    \theutterance \stepcounter{utterance}  
    & & \multicolumn{4}{p{0.6\linewidth}}{
        \cellcolor[rgb]{0.9,0.9,0.9}{
            \makecell[{{p{\linewidth}}}]{
                \texttt{\tiny{[GM$\rangle$P2]}}
                \texttt{Słowa tajfun, tornado, huragan zostały odgadnięte prawidłowo. Twoja tura się tam skończyła.} \\
\texttt{Teraz podaj kolejną listę oddzieloną przecinkami, zawierającą co najmniej 1 i maksymalnie 3 słów z poniższej listy słów, które najlepiej odnoszą się lub są najbardziej związane ze słowem 'przestępstwo'. Pamiętaj, aby rozpocząć odpowiedź od ZGADNIĘCIE: . Uwaga: niektóre słowa zostały usunięte z listy w porównaniu do poprzednich próśb.} \\
\\ 
\texttt{mechanik, rozbój, deszcz, krawat, kardynał, strzyżyk, róża, stopa, pistolet, bąk, prawnik, pelargonia, cymbał, kilometr, pasek, bluza, księgowy, oszustwo, jard, przekroczenie prędkości, koszula} \\
            }
        }
    }
    & \\ \\

    \theutterance \stepcounter{utterance}  
    & & & \multicolumn{4}{p{0.6\linewidth}}{
        \cellcolor[rgb]{0.9,0.9,0.9}{
            \makecell[{{p{\linewidth}}}]{
                \texttt{\tiny{[GM$\langle$P2]}}
                \texttt{ZGADNIĘCIE: rozbój, oszustwo, przekroczenie prędkości} \\
            }
        }
    }
     \\ \\

    \theutterance \stepcounter{utterance}  
    & & & \multicolumn{2}{p{0.3\linewidth}}{
        \cellcolor[rgb]{0.9,0.9,0.9}{
            \makecell[{{p{\linewidth}}}]{
                \texttt{\tiny{[GM$|$GM]}}
                ['rozbój', 'oszustwo', 'przekroczenie prędkości']
            }
        }
    }
    & & \\ \\

    \theutterance \stepcounter{utterance}  
    & & & \multicolumn{2}{p{0.3\linewidth}}{
        \cellcolor[rgb]{0.9,0.9,0.9}{
            \makecell[{{p{\linewidth}}}]{
                \texttt{\tiny{[GM$|$GM]}}
                {'word': 'rozbój', 'assignment': 'team'}
            }
        }
    }
    & & \\ \\

    \theutterance \stepcounter{utterance}  
    & & & \multicolumn{2}{p{0.3\linewidth}}{
        \cellcolor[rgb]{0.9,0.9,0.9}{
            \makecell[{{p{\linewidth}}}]{
                \texttt{\tiny{[GM$|$GM]}}
                {'word': 'rozbój', 'assignment': 'team'}
            }
        }
    }
    & & \\ \\

    \theutterance \stepcounter{utterance}  
    & & & \multicolumn{2}{p{0.3\linewidth}}{
        \cellcolor[rgb]{0.9,0.9,0.9}{
            \makecell[{{p{\linewidth}}}]{
                \texttt{\tiny{[GM$|$GM]}}
                {'word': 'oszustwo', 'assignment': 'team'}
            }
        }
    }
    & & \\ \\

    \theutterance \stepcounter{utterance}  
    & & & \multicolumn{2}{p{0.3\linewidth}}{
        \cellcolor[rgb]{0.9,0.9,0.9}{
            \makecell[{{p{\linewidth}}}]{
                \texttt{\tiny{[GM$|$GM]}}
                {'word': 'oszustwo', 'assignment': 'team'}
            }
        }
    }
    & & \\ \\

    \theutterance \stepcounter{utterance}  
    & & & \multicolumn{2}{p{0.3\linewidth}}{
        \cellcolor[rgb]{0.9,0.9,0.9}{
            \makecell[{{p{\linewidth}}}]{
                \texttt{\tiny{[GM$|$GM]}}
                {'word': 'przekroczenie prędkości', 'assignment': 'team'}
            }
        }
    }
    & & \\ \\

    \theutterance \stepcounter{utterance}  
    & & & \multicolumn{2}{p{0.3\linewidth}}{
        \cellcolor[rgb]{0.9,0.9,0.9}{
            \makecell[{{p{\linewidth}}}]{
                \texttt{\tiny{[GM$|$GM]}}
                {'word': 'przekroczenie prędkości', 'assignment': 'team'}
            }
        }
    }
    & & \\ \\

    \theutterance \stepcounter{utterance}  
    & & & \multicolumn{2}{p{0.3\linewidth}}{
        \cellcolor[rgb]{0.9,0.9,0.9}{
            \makecell[{{p{\linewidth}}}]{
                \texttt{\tiny{[GM$|$GM]}}
                {'word': 'prawnik', 'assignment': 'opponent'}
            }
        }
    }
    & & \\ \\

    \theutterance \stepcounter{utterance}  
    & & & \multicolumn{2}{p{0.3\linewidth}}{
        \cellcolor[rgb]{0.9,0.9,0.9}{
            \makecell[{{p{\linewidth}}}]{
                \texttt{\tiny{[GM$|$GM]}}
                {'hidden': {'team': ['strzyżyk', 'deszcz', 'kardynał'], 'innocent': ['stopa', 'koszula', 'jard', 'krawat', 'pasek', 'kilometr', 'bluza'], 'opponent': ['róża', 'księgowy', 'pistolet', 'bąk', 'pelargonia', 'mechanik'], 'assassin': ['cymbał']}, 'revealed': {'team': {'team': ['tajfun', 'tornado', 'huragan', 'rozbój', 'oszustwo', 'przekroczenie prędkości'], 'innocent': [], 'opponent': [], 'assassin': []}, 'opponent': {'team': [], 'innocent': [], 'opponent': ['nagietek', 'prawnik'], 'assassin': []}}}
            }
        }
    }
    & & \\ \\

    \theutterance \stepcounter{utterance}  
    & & \multicolumn{4}{p{0.6\linewidth}}{
        \cellcolor[rgb]{0.9,0.9,0.9}{
            \makecell[{{p{\linewidth}}}]{
                \texttt{\tiny{[P1$\langle$GM]}}
                \texttt{Słowa rozbój, oszustwo, przekroczenie prędkości zostały odgadnięte prawidłowo. Tura twojego kolegi z drużyny się tam skończyła.} \\
\texttt{Teraz podaj kolejną wskazówkę odnoszącą się do niektórych z pozostałych słów twojego zespołu oraz listę powiązanych słów docelowych. Pamiętaj, aby rozpocząć wskazówkę od 'WSKAZÓWKA: ', wstawić nową linię i rozpocząć listę słów docelowych oddzielonych przecinkami od 'CELE: '. Uwaga: niektóre słowa zostały usunięte z list w porównaniu do poprzednich żądań.} \\
\\ 
\texttt{Twoje pozostałe słowa zespołu to: strzyżyk, deszcz, kardynał.} \\
\texttt{Pozostałe słowa twojego przeciwnika to: róża, księgowy, pistolet, bąk, pelargonia, prawnik, mechanik.} \\
\texttt{Pozostałe słowa rozpraszające to: stopa, koszula, jard, krawat, pasek, kilometr, bluza.} \\
\texttt{Pozostałe słowa zabójcy to: cymbał.} \\
            }
        }
    }
    & \\ \\

    \theutterance \stepcounter{utterance}  
    & \multicolumn{4}{p{0.6\linewidth}}{
        \cellcolor[rgb]{0.9,0.9,0.9}{
            \makecell[{{p{\linewidth}}}]{
                \texttt{\tiny{[P1$\rangle$GM]}}
                \texttt{WSKAZÓWKA: ptak} \\
\texttt{CELE: strzyżyk, kardynał} \\
            }
        }
    }
    & & \\ \\

    \theutterance \stepcounter{utterance}  
    & & & \multicolumn{2}{p{0.3\linewidth}}{
        \cellcolor[rgb]{0.9,0.9,0.9}{
            \makecell[{{p{\linewidth}}}]{
                \texttt{\tiny{[GM$|$GM]}}
                \texttt{ptak} \\
            }
        }
    }
    & & \\ \\

    \theutterance \stepcounter{utterance}  
    & & & \multicolumn{2}{p{0.3\linewidth}}{
        \cellcolor[rgb]{0.9,0.9,0.9}{
            \makecell[{{p{\linewidth}}}]{
                \texttt{\tiny{[GM$|$GM]}}
                ['strzyżyk', 'kardynał']
            }
        }
    }
    & & \\ \\

    \theutterance \stepcounter{utterance}  
    & & & \multicolumn{2}{p{0.3\linewidth}}{
        \cellcolor[rgb]{0.9,0.9,0.9}{
            \makecell[{{p{\linewidth}}}]{
                \texttt{\tiny{[GM$|$GM]}}
                {'word': 'strzyżyk', 'assignment': 'team'}
            }
        }
    }
    & & \\ \\

    \theutterance \stepcounter{utterance}  
    & & & \multicolumn{2}{p{0.3\linewidth}}{
        \cellcolor[rgb]{0.9,0.9,0.9}{
            \makecell[{{p{\linewidth}}}]{
                \texttt{\tiny{[GM$|$GM]}}
                {'word': 'kardynał', 'assignment': 'team'}
            }
        }
    }
    & & \\ \\

    \theutterance \stepcounter{utterance}  
    & & \multicolumn{4}{p{0.6\linewidth}}{
        \cellcolor[rgb]{0.9,0.9,0.9}{
            \makecell[{{p{\linewidth}}}]{
                \texttt{\tiny{[GM$\rangle$P2]}}
                \texttt{Słowa rozbój, oszustwo, przekroczenie prędkości zostały odgadnięte prawidłowo. Twoja tura się tam skończyła.} \\
\texttt{Teraz podaj kolejną listę oddzieloną przecinkami, zawierającą co najmniej 1 i maksymalnie 2 słów z poniższej listy słów, które najlepiej odnoszą się lub są najbardziej związane ze słowem 'ptak'. Pamiętaj, aby rozpocząć odpowiedź od ZGADNIĘCIE: . Uwaga: niektóre słowa zostały usunięte z listy w porównaniu do poprzednich próśb.} \\
\\ 
\texttt{mechanik, deszcz, krawat, kardynał, strzyżyk, róża, stopa, pistolet, bąk, pelargonia, cymbał, kilometr, pasek, bluza, księgowy, jard, koszula} \\
            }
        }
    }
    & \\ \\

    \theutterance \stepcounter{utterance}  
    & & & \multicolumn{4}{p{0.6\linewidth}}{
        \cellcolor[rgb]{0.9,0.9,0.9}{
            \makecell[{{p{\linewidth}}}]{
                \texttt{\tiny{[GM$\langle$P2]}}
                \texttt{ZGADNIĘCIE: strzyżyk, kardynał} \\
            }
        }
    }
     \\ \\

    \theutterance \stepcounter{utterance}  
    & & & \multicolumn{2}{p{0.3\linewidth}}{
        \cellcolor[rgb]{0.9,0.9,0.9}{
            \makecell[{{p{\linewidth}}}]{
                \texttt{\tiny{[GM$|$GM]}}
                ['strzyżyk', 'kardynał']
            }
        }
    }
    & & \\ \\

    \theutterance \stepcounter{utterance}  
    & & & \multicolumn{2}{p{0.3\linewidth}}{
        \cellcolor[rgb]{0.9,0.9,0.9}{
            \makecell[{{p{\linewidth}}}]{
                \texttt{\tiny{[GM$|$GM]}}
                {'word': 'strzyżyk', 'assignment': 'team'}
            }
        }
    }
    & & \\ \\

    \theutterance \stepcounter{utterance}  
    & & & \multicolumn{2}{p{0.3\linewidth}}{
        \cellcolor[rgb]{0.9,0.9,0.9}{
            \makecell[{{p{\linewidth}}}]{
                \texttt{\tiny{[GM$|$GM]}}
                {'word': 'strzyżyk', 'assignment': 'team'}
            }
        }
    }
    & & \\ \\

    \theutterance \stepcounter{utterance}  
    & & & \multicolumn{2}{p{0.3\linewidth}}{
        \cellcolor[rgb]{0.9,0.9,0.9}{
            \makecell[{{p{\linewidth}}}]{
                \texttt{\tiny{[GM$|$GM]}}
                {'word': 'kardynał', 'assignment': 'team'}
            }
        }
    }
    & & \\ \\

    \theutterance \stepcounter{utterance}  
    & & & \multicolumn{2}{p{0.3\linewidth}}{
        \cellcolor[rgb]{0.9,0.9,0.9}{
            \makecell[{{p{\linewidth}}}]{
                \texttt{\tiny{[GM$|$GM]}}
                {'word': 'kardynał', 'assignment': 'team'}
            }
        }
    }
    & & \\ \\

    \theutterance \stepcounter{utterance}  
    & & & \multicolumn{2}{p{0.3\linewidth}}{
        \cellcolor[rgb]{0.9,0.9,0.9}{
            \makecell[{{p{\linewidth}}}]{
                \texttt{\tiny{[GM$|$GM]}}
                {'word': 'bąk', 'assignment': 'opponent'}
            }
        }
    }
    & & \\ \\

    \theutterance \stepcounter{utterance}  
    & & & \multicolumn{2}{p{0.3\linewidth}}{
        \cellcolor[rgb]{0.9,0.9,0.9}{
            \makecell[{{p{\linewidth}}}]{
                \texttt{\tiny{[GM$|$GM]}}
                {'hidden': {'team': ['deszcz'], 'innocent': ['stopa', 'koszula', 'jard', 'krawat', 'pasek', 'kilometr', 'bluza'], 'opponent': ['róża', 'księgowy', 'pistolet', 'pelargonia', 'mechanik'], 'assassin': ['cymbał']}, 'revealed': {'team': {'team': ['tajfun', 'tornado', 'huragan', 'rozbój', 'oszustwo', 'przekroczenie prędkości', 'strzyżyk', 'kardynał'], 'innocent': [], 'opponent': [], 'assassin': []}, 'opponent': {'team': [], 'innocent': [], 'opponent': ['nagietek', 'prawnik', 'bąk'], 'assassin': []}}}
            }
        }
    }
    & & \\ \\

    \theutterance \stepcounter{utterance}  
    & & \multicolumn{4}{p{0.6\linewidth}}{
        \cellcolor[rgb]{0.9,0.9,0.9}{
            \makecell[{{p{\linewidth}}}]{
                \texttt{\tiny{[P1$\langle$GM]}}
                \texttt{Słowa strzyżyk, kardynał zostały odgadnięte prawidłowo. Tura twojego kolegi z drużyny się tam skończyła.} \\
\texttt{Teraz podaj kolejną wskazówkę odnoszącą się do niektórych z pozostałych słów twojego zespołu oraz listę powiązanych słów docelowych. Pamiętaj, aby rozpocząć wskazówkę od 'WSKAZÓWKA: ', wstawić nową linię i rozpocząć listę słów docelowych oddzielonych przecinkami od 'CELE: '. Uwaga: niektóre słowa zostały usunięte z list w porównaniu do poprzednich żądań.} \\
\\ 
\texttt{Twoje pozostałe słowa zespołu to: deszcz.} \\
\texttt{Pozostałe słowa twojego przeciwnika to: róża, księgowy, pistolet, bąk, pelargonia, mechanik.} \\
\texttt{Pozostałe słowa rozpraszające to: stopa, koszula, jard, krawat, pasek, kilometr, bluza.} \\
\texttt{Pozostałe słowa zabójcy to: cymbał.} \\
            }
        }
    }
    & \\ \\

    \theutterance \stepcounter{utterance}  
    & \multicolumn{4}{p{0.6\linewidth}}{
        \cellcolor[rgb]{0.9,0.9,0.9}{
            \makecell[{{p{\linewidth}}}]{
                \texttt{\tiny{[P1$\rangle$GM]}}
                \texttt{WSKAZÓWKA: opad} \\
\texttt{CELE: deszcz} \\
            }
        }
    }
    & & \\ \\

    \theutterance \stepcounter{utterance}  
    & & & \multicolumn{2}{p{0.3\linewidth}}{
        \cellcolor[rgb]{0.9,0.9,0.9}{
            \makecell[{{p{\linewidth}}}]{
                \texttt{\tiny{[GM$|$GM]}}
                \texttt{opad} \\
            }
        }
    }
    & & \\ \\

    \theutterance \stepcounter{utterance}  
    & & & \multicolumn{2}{p{0.3\linewidth}}{
        \cellcolor[rgb]{0.9,0.9,0.9}{
            \makecell[{{p{\linewidth}}}]{
                \texttt{\tiny{[GM$|$GM]}}
                ['deszcz']
            }
        }
    }
    & & \\ \\

    \theutterance \stepcounter{utterance}  
    & & & \multicolumn{2}{p{0.3\linewidth}}{
        \cellcolor[rgb]{0.9,0.9,0.9}{
            \makecell[{{p{\linewidth}}}]{
                \texttt{\tiny{[GM$|$GM]}}
                {'word': 'deszcz', 'assignment': 'team'}
            }
        }
    }
    & & \\ \\

    \theutterance \stepcounter{utterance}  
    & & \multicolumn{4}{p{0.6\linewidth}}{
        \cellcolor[rgb]{0.9,0.9,0.9}{
            \makecell[{{p{\linewidth}}}]{
                \texttt{\tiny{[GM$\rangle$P2]}}
                \texttt{Słowa strzyżyk, kardynał zostały odgadnięte prawidłowo. Twoja tura się tam skończyła.} \\
\texttt{Teraz podaj kolejną listę oddzieloną przecinkami, zawierającą co najmniej 1 i maksymalnie 1 słów z poniższej listy słów, które najlepiej odnoszą się lub są najbardziej związane ze słowem 'opad'. Pamiętaj, aby rozpocząć odpowiedź od ZGADNIĘCIE: . Uwaga: niektóre słowa zostały usunięte z listy w porównaniu do poprzednich próśb.} \\
\\ 
\texttt{mechanik, deszcz, krawat, róża, stopa, pistolet, pelargonia, cymbał, kilometr, pasek, bluza, księgowy, jard, koszula} \\
            }
        }
    }
    & \\ \\

    \theutterance \stepcounter{utterance}  
    & & & \multicolumn{4}{p{0.6\linewidth}}{
        \cellcolor[rgb]{0.9,0.9,0.9}{
            \makecell[{{p{\linewidth}}}]{
                \texttt{\tiny{[GM$\langle$P2]}}
                \texttt{ZGADNIĘCIE: deszcz} \\
            }
        }
    }
     \\ \\

    \theutterance \stepcounter{utterance}  
    & & & \multicolumn{2}{p{0.3\linewidth}}{
        \cellcolor[rgb]{0.9,0.9,0.9}{
            \makecell[{{p{\linewidth}}}]{
                \texttt{\tiny{[GM$|$GM]}}
                ['deszcz']
            }
        }
    }
    & & \\ \\

    \theutterance \stepcounter{utterance}  
    & & & \multicolumn{2}{p{0.3\linewidth}}{
        \cellcolor[rgb]{0.9,0.9,0.9}{
            \makecell[{{p{\linewidth}}}]{
                \texttt{\tiny{[GM$|$GM]}}
                {'word': 'deszcz', 'assignment': 'team'}
            }
        }
    }
    & & \\ \\

    \theutterance \stepcounter{utterance}  
    & & & \multicolumn{2}{p{0.3\linewidth}}{
        \cellcolor[rgb]{0.9,0.9,0.9}{
            \makecell[{{p{\linewidth}}}]{
                \texttt{\tiny{[GM$|$GM]}}
                {'word': 'deszcz', 'assignment': 'team'}
            }
        }
    }
    & & \\ \\

    \theutterance \stepcounter{utterance}  
    & & & \multicolumn{2}{p{0.3\linewidth}}{
        \cellcolor[rgb]{0.9,0.9,0.9}{
            \makecell[{{p{\linewidth}}}]{
                \texttt{\tiny{[GM$|$GM]}}
                \texttt{team has won} \\
            }
        }
    }
    & & \\ \\

\end{supertabular}
}

\end{document}
